
% %%%%%%%%%%%%%%%%%%%%%%%%%%%%%%%%%%%%%%%%%%%%%%%%%%%%%%%%%%%%%%%%%%%%%%%%%%%%%%%
% % RpgDropCap
% %%%%%%%%%%%%%%%%%%%%%%%%%%%%%%%%%%%%%%%%%%%%%%%%%%%%%%%%%%%%%%%%%%%%%%%%%%%%%%%
\ExplSyntaxOn
\renewcommand{\LettrineFontHook} { \RpgFontDropCap} %feed the desired font into the lettrine pipeline
\renewcommand{\LettrineTextFont}{\RpgFontDropCapAfter}

% Usage: rpgDropCapLine[<lettrine options>]{<first letter>}{<small caps line>}
%    It takes trial and error to get the 2nd argument to align with
%    the linebreak. See the lettrine package for options.
\NewDocumentCommand{\RpgDropCap}{ O{} m m }
{
	\lettrine[
		lines   = \__rpg_DropCapHeight,
		depth   = 0,
		findent = \RpgDropCapSeparation,
		nindent = 0pt,
		#1
	]
	{#2}{#3}
}


%%%%%%%%%%%%%%%%%%%%%%%%%%%%%%
% Global lettrine parameters
%%%%%%%%%%%%%%%%%%%%%%%%%%%%%%

% horizontal length after lettrine (modify with \setlength) 
\newlength\RpgDropCapSeparation
\setlength{\RpgDropCapSeparation}{\dim_use:N \__rpg_SpaceDim}

% integer count of lines taken up by drop cap
\int_new:N \__rpg_DropCapHeight

%% user interface for setting the line count
\NewDocumentCommand{\RpgThemeDropCapHeight}{ m }
{
	\int_set:Nn\__rpg_DropCapHeight{#1}
}
\RpgThemeDropCapHeight{4}
